\section{August 26, 2026}

\subsection{Further properties of fields}

The preceding lecture showed that $a\cdot 0=0$ for every element $a$ of
a field. We now obtain the zero-product property.

\begin{proposition}[Zero-product property]
  Let $a,b\in\bb{F}$. If $ab=0$, then either $a=0$ or $b=0$.
\end{proposition}

\begin{proof}
  Suppose that $ab=0$ and $a\neq0$. Since $\bb{F}$ is a field,
  $a^{-1}$ exists. Therefore
  \[
    0=a^{-1}0=a^{-1}(ab)=(a^{-1}a)b=1b=b.
  \]
  Thus $a=0$ or $b=0$.
\end{proof}

\subsection{Characteristic of a field}

For a positive integer $n$ and $a\in\bb{F}$, write
\[
  na=\underbrace{a+a+\cdots+a}_{n\text{ terms}}.
\]

\begin{definition}[Characteristic]
  Let $\bb{F}$ be a field. If $n1=0$ for some positive integer $n$, then
  the least such $n$ is called the \emph{characteristic} of $\bb{F}$,
  denoted $\operatorname{char}(\bb{F})$. If $n1\neq0$ for every positive
  integer $n$, then $\bb{F}$ has characteristic zero.
\end{definition}

The fields $\bb{Q}$, $\bb{R}$, and $\bb{C}$ have characteristic zero.
For every prime $p$, the field $\bb{F}_p$ has characteristic $p$. The
rational function field $\bb{F}_p(t)$ is an infinite field of
characteristic $p$.

\begin{proposition}
  If $\bb{F}$ has positive characteristic $r$, then $r$ is prime.
\end{proposition}

\begin{proof}
  We have $r\neq1$, since $1\neq0$. Suppose that $r$ is composite, so
  \[
    r=st
  \]
  for integers $s,t$ satisfying $1<s,t<r$. Distributivity gives
  \[
    (s1)(t1)=(st)1=r1=0.
  \]
  By the zero-product property, either $s1=0$ or $t1=0$. Both
  alternatives contradict the minimality of $r$. Hence $r$ is prime.
\end{proof}

\begin{proposition}
  If $\operatorname{char}(\bb{F})=p>0$, then
  \[
    pa=0
  \]
  for every $a\in\bb{F}$.
\end{proposition}

\begin{proof}
  Since $p1=0$, associativity and distributivity imply
  \[
    pa=p(1a)=(p1)a=0a=0.
  \]
\end{proof}

\begin{proposition}
  Every finite field has positive characteristic.
\end{proposition}

\begin{proof}
  Let $\bb{F}$ be finite. The sequence
  \[
    1,\quad 2\cdot1,\quad 3\cdot1,\quad\ldots
  \]
  cannot consist of distinct elements. Hence there are positive integers
  $m>n$ such that
  \[
    m1=n1.
  \]
  Adding the additive inverse of $n1$ to both sides gives
  \[
    (m-n)1=0.
  \]
  Thus $\bb{F}$ has positive characteristic.
\end{proof}

\subsection{Algebraically closed fields}

\begin{definition}[Algebraically closed field]
  A field $\bb{F}$ is \emph{algebraically closed} if every nonconstant
  polynomial $f\in\bb{F}[x]$ has a root in $\bb{F}$.
\end{definition}

The field $\bb{Q}$ is not algebraically closed, because $x^2-2$ has no
root in $\bb{Q}$. The field $\bb{R}$ is not algebraically closed, because
$x^2+1$ has no real root. By the fundamental theorem of algebra,
$\bb{C}$ is algebraically closed.

\begin{proposition}
  No finite field is algebraically closed.
\end{proposition}

\begin{proof}
  Let
  \[
    \bb{F}=\{a_1,\ldots,a_q\}
  \]
  be finite. Consider the nonconstant polynomial
  \[
    f(x)=1+\prod_{i=1}^q(x-a_i).
  \]
  For each $i=1,\ldots,q$,
  \[
    f(a_i)=1\neq0.
  \]
  Thus $f$ has no root in $\bb{F}$, so $\bb{F}$ is not algebraically
  closed.
\end{proof}

\subsection{Vectors and vector spaces}

The spaces $\bb{R}^2$, $\bb{R}^3$, and, more generally,
\[
  \bb{R}^n
  =\left\{
    \begin{bmatrix}
      x_1\\ \vdots\\ x_n
    \end{bmatrix}
    :x_1,\ldots,x_n\in\bb{R}
  \right\}
\]
are the motivating examples of vector spaces.

\begin{definition}[Vector space]
  Let $\bb{F}$ be a field. A \emph{vector space over $\bb{F}$} is a set
  $V$ equipped with two operations,
  \[
    V\times V\longrightarrow V,
    \qquad
    (u,v)\longmapsto u+v,
  \]
  and
  \[
    \bb{F}\times V\longrightarrow V,
    \qquad
    (\alpha,v)\longmapsto\alpha v,
  \]
  satisfying the following axioms for all $u,v,w\in V$ and
  $\alpha,\beta\in\bb{F}$:
  \begin{enumerate}[label=(V\arabic*)]
    \item $u+(v+w)=(u+v)+w$;
    \item $u+v=v+u$;
    \item there is an element $0\in V$ such that $v+0=0+v=v$;
    \item for every $v\in V$, there is an element $-v\in V$ such that
      $v+(-v)=(-v)+v=0$;
    \item $1v=v$;
    \item $(\alpha\beta)v=\alpha(\beta v)$;
    \item $\alpha(u+v)=\alpha u+\alpha v$;
    \item $(\alpha+\beta)v=\alpha v+\beta v$.
  \end{enumerate}
\end{definition}

The elements of $V$ are called \emph{vectors}, and the elements of
$\bb{F}$ are called \emph{scalars}.
